%%%%%%%%%%%%%%%%%%%%%%%%%%%%%%%%%%%%%%%%%%%%%%%%%%%%%%%%%%%%%%%%%%%%%%%%%%%%%%%
\section{Related Work}\label{sec:related}
%%%%%%%%%%%%%%%%%%%%%%%%%%%%%%%%%%%%%%%%%%%%%%%%%%%%%%%%%%%%%%%%%%%%%%%%%%%%%%%

In addition to the ZEBRA, which we have described previously, there have been several previous attempts at developing wearable-authentication technology.  Two of note include the Bionym Nymi \cite{goode2014bring} and the Intel Authentication Bracelet \cite{idf2015intel}.  The Bionym Nymi is an authentication wristband that broadcasts a digitally signed authentication signal derived from a user's heartbeat to nearby devices using BLE~\cite{BLEoverview}.  The Intel Authentication Bracelet requires a user to log in to a system with a standard password.  A credential is then transmitted to the bracelet using BLE.  This credential is then broadcast to nearby bracelet-enabled devices in order to allow password-less login.

\subsubsection{Limitations of Existing Work}
Existing authentication wearables use wireless communication technology (typically BLE) to broadcast their authentication credentials.  Thus the devices are likely vulnerable to ghost-and-leech attacks \cite{czeskis2008rfids}.  Ghost-and-leech attacks occur when an attacker uses a more powerful radio transmitter than the transmitter found on a wireless device in order to capture and rebroadcast the wireless signal in order to fool a target into believing that the wireless device is in closer proximity to the target than it actually is.

%Related work can be divided into two key areas: authentication and deauthentication.  Authentication-based research looks at using wristbands or other wearables to help authenticate a user to a computing resource.  Deauthentication-based research focuses on wearables that do the opposite, helping to deauthenticate a user from a computer resource when the user is done using the resource.  Some wearables are able to both authenticate users to and deauthenticate users from computing resources \cite{} while other wearables can only be used for authentication or for deauthentication \cite{}.

%\subsection{Authentication}
%There have been several previous attempts at developing wearable-authentication technology.  Some of these attempts have taken drastically different approaches than kbid while others share some similarities.  However, none of the existing technology has considered using skin conductance as the transfer vehicle for the authentication token in order to provide robustness against repeater attacks.

%\subsubsection{Bionym Nymi}
%The Bionym Nymi \cite{} is an authentication wristband that broadcasts a digitally signed authentication signal derived from a user's heartbeat to nearby devices using Bluetooth Low Energy.  "The Nymi's ECG recognition algorithms observe the shape of the ECG waveform, extracting unique and consistent features that are a result of a user's physiology. The Nymi is robust against a wide variety of health issues." \cite{}.

%\subsubsection{Intel Authentication Bracelet Prototype}
%Intel, in partnership with Fossil, has announced a prototype for an authentication bracelet based on its Curie platform \cite{}.  The bracelet is designed to automatically unlock a device when a given user is nearby.  When the bracelet is first put on, the user must log in to a bracelet-enabled device with a standard password.  However, subsequent logins to bracelet-enabled devices do not require a password and the bracelet will continue to broadcast an authentication token until the bracelet has been taken off, causing the device's volatile storage to lose the authentication credential.  The bracelet uses Bluetooth Low Energy to broadcast its credential to nearby bracelet-enabled devices.

%\subsubsection{Limitations of Existing Work}
%Existing authentication wearables use wireless communication technology (typically Bluetooth LE)to broadcast their authentication credentials.  Thus the devices are likely vulnerable to ghost-and-leech attacks.  Ghost-and-leech attacks occur when an attacker uses a more powerful radio transmitter than the transmitter found on a wireless device in order to capture and rebroadcast the wireless signal.  

%\par
%In our context, an attacker could capture and rebroadcast the authentication signal off of an authentication wearable and thus could log in to a computer terminal even when an authorized user is not near the terminal.  For example, in a hospital, a doctor might wear an authentication bracelet which continuously broadcasts a changing login credential using Bluetooth Low Energy.  An attacker might place a receiver next to a target doctor's office.  The receiver could then capture the signal and send it to a portable rebroadcaster which the attacker could use to unlock a computer terminal far from the doctor's office.  Due to the wireless nature of these types of authentication wearables, these types of attacks can be extremely difficult to protect against.

%\subsection{Deauthentication}
%There are numerous wearables that are able to successfully deauthenticate a user when the physical connection between the wearable and the computer is broken \cite{}.  Most of these techniques form a circuit that can be physically broken when the wearable is removed \cite{} but some more advanced techniques such as Apple's "Wrist Detect" \cite{} can determine if a wearable is in constant contact with skin.

%\subsubsection{ZEBRA: Zero-Effort Bilateral Recurring Authentication}
%ZEBRA \cite{} was developed to solve a very different problem than the authentication bracelets described above.  Instead of authenticating a user to a computer terminal, ZEBRA is focused on detecting when a user is done using the computer terminal and deauthenticating the user from the terminal even if the user forgets to do so.  ZEBRA is prototyped as a bracelet containing an accelerometer and a variety of sensors.  A user puts on the bracelet and begins using a computer terminal as he or she normally would.  ZEBRA dynamically builds a profile of the user's hand movements in a short amount of time and then enforces this profile for the duration of the user's computing session by comparing the hand movements observed with the wrist band to data captured from keyboard and mouse movements.
